\chapter{Методика}
\label{ch:methods}

    Молекулярно-кинетическая теория дает следующую оценку для коэффициента теплопроводности газов:

    \begin{align}
        k &\sim \lambda \overline{\nu} \cdot n c_V \label{k},
    \end{align}
    где $\lambda$ - длина свободного пробега молекул газа, $\overline{\nu} = \sqrt{\frac{8k_{\text{Б}}T}{\uppi m}}$ — средняя скорость их теплового движения, $n$ — концентрация (объёмная плотность) газа, $C_V = \frac{i}{2}k_{\text{Б}}$ — его теплоёмкость при постоянном объёме в расчёте на одну молекулу ($i$ — эффективное число степеней свободы молекулы).

    Длина свободного пробега может быть оценена как $\lambda = 1/n\upsigma$, где $\upsigma$ — эффективное сечение столкновений молекул друг с другом. Тогда из (2) видно, что коэффициент теплопроводности газа не зависит от плотности газа и определяется только его температурой. В простейшей модели твёрдых шариков $\upsigma = const$, и коэффициент теплопроводности пропорционален корню абсолютной температуры: $k \sim \overline{\nu}/\upsigma \sim \sqrt{T}$. На практике эффективное сечение $\upsigma(T)$ следует считать медленно убывающей функцией $T$. Поэтому необходимо измерять коэффициент теплопроводности экспериментально.

    %что-то про саму методику...
    Экспериментально определить коэффициент теплопроводности можно, используя закон Фурье о плотности потока энергии.
        

   